%%%%%%%%%%%%%%%%%%%%%%%%%%%%%%%%%%%%%%%%%%%%%%%%%%%%
%%%     Language Science Press Master File       %%%
%%%              edited volumes                  %%%
%%%         follow the instructions below        %%%
%%%%%%%%%%%%%%%%%%%%%%%%%%%%%%%%%%%%%%%%%%%%%%%%%%%%

% Everything following a % is ignored
% Some lines start with %. Remove the % to include them

\documentclass[output=book,collection%,chinesefont
%             ,arabicfont
%             ,colorlinks, citecolor=brown
            % ,draft
%            ,draftmode
		  ]{langscibook}
%%%%%%%%%%%%%%%%%%%%%%%%%%%%%%%%%%%%%%%%%%%%%%%%%%%%

% put all additional commands you need in the
% following files. If you do not know what this might
% mean, you can safely ignore this section

\title{Advancing interdisciplinarity in empirical translation research}
% \subtitle{Add subtitle here if it exists}
\BackBody{Thanks to its broad scope of topics and linguistic diversity, this volume reflects the vast range of research sub-fields in contemporary translation and interpreting studies, including specialized translation, machine translation, and audio-visual translation. At the same time, the chapters illustrate the intrinsic interdisciplinary foundations of translation studies, highlighting its connection to specific linguistic sub-fields and applied branches. Among the linguistic fields, the strongest ties here are with contrastive linguistics and corpus linguistics, while its connection to practical applications is evident particularly in its focus on developing effective translator training methods.
}
\author{Tomáš Svoboda and Olga Nádvorníková and Věra Kloudová}

\renewcommand{\lsISBNdigital}{978-3-96110-595-3}
\renewcommand{\lsISBNhardcover}{978-3-98554-212-3}
\BookDOI{10.5281/zenodo.22148115}
\typesetter{Sebastian Nordhoff}
\proofreader{Matthew Korte}
% \lsCoverTitleSizes{51.5pt}{17pt}% Font setting for the title page


\renewcommand{\lsSeries}{tmnlp}
\renewcommand{\lsID}{574}
\renewcommand{\lsSeriesNumber}{27}

\usepackage{tabularx,multicol}
\usepackage{url}
\urlstyle{same}

 
\usepackage{langsci-optional}
\usepackage{langsci-lgr}
\usepackage{langsci-gb4e}

\usepackage{paralist}
\usepackage{enumitem}
\usepackage{cprotect}
\usepackage{fancyvrb}

\let\clipbox\relax
\usepackage[export]{adjustbox}
\usepackage{array}
\usepackage[normalem]{ulem}
\usepackage{langsci-branding}
\usepackage{langsci-bidi}
\usepackage[linguistics]{forest}

%% hyphenation points for line breaks
%% Normally, automatic hyphenation in LaTeX is very good
%% If a word is mis-hyphenated, add it to this file
%%
%% add information to TeX file before \begin{document} with:
%% %% hyphenation points for line breaks
%% Normally, automatic hyphenation in LaTeX is very good
%% If a word is mis-hyphenated, add it to this file
%%
%% add information to TeX file before \begin{document} with:
%% %% hyphenation points for line breaks
%% Normally, automatic hyphenation in LaTeX is very good
%% If a word is mis-hyphenated, add it to this file
%%
%% add information to TeX file before \begin{document} with:
%% \include{localhyphenation}
\hyphenation{
    par-a-digm
}

\hyphenation{
    par-a-digm
}

\hyphenation{
    par-a-digm
}


\makeatletter
\let\thetitle\@title
\let\theauthor\@author
\makeatother

\newcommand{\togglepaper}[1][0]{
   \bibliography{../localbibliography}
   \papernote{\scriptsize\normalfont
     \theauthor.
     \titleTemp.
     To appear in:
     E. Di Tor \& Herr Rausgeberin (ed.).
     Booktitle in localcommands.tex.
     Berlin: Language Science Press. [preliminary page numbering]
   }
   \pagenumbering{roman}
   \setcounter{chapter}{#1}
   \addtocounter{chapter}{-1}
}

\newbool{bookcompile}
\booltrue{bookcompile}
\newcommand{\bookorchapter}[2]{\ifbool{bookcompile}{#1}{#2}}

\newcommand{\qual}[1]{`#1'}
\newcommand{\bv}[1]{(\ref{#1})}

\newcommand{\ip}{ImPost}
\newcommand{\ipadv}{ImPost-adverbial}
\newcommand{\pradv}{ProxPost-adverbial}
\newcommand{\chisq}{Chi-squared}


% for rotated columns
\newcolumntype{T}[2]{%
    >{\adjustbox{angle=#1,lap=\width-(#2)}\bgroup}%
    l%
    <{\egroup}%
}
\newcommand*\rot{\multicolumn{1}{T{45}{1em}}}%


\VerbatimFootnotes


% \newcommand{\keywords}[1]{\textbf{#1}}


%  \newcommand{\textarabic}[1]{\RL{\arabicfont #1}}

\addbibresource{localbibliography.bib}
\addbibresource{new_translatology.bib}
\addbibresource{gast.bib}

\newfontfamily\cn
	  [
	    Scale=MatchLowercase,
	    BoldFont=SourceHanSerif-Bold.otf
      ]
      {SourceHanSerif-Regular.otf}
    \AdditionalFontImprint{Source Han Serif ZH}
	\XeTeXlinebreaklocale 'zh'
	\XeTeXlinebreakskip = 0pt plus 1pt


\newfontfamily\arabicfont[Script=Arabic,ItalicFont=*,Scale=1.4]{ScheherazadeRegOT_Jazm.ttf}
    \AdditionalFontImprint{Arabtype}
    \usepackage{langsci-bidi}


%%%%%%%%%%%%%%%%%%%%%%%%%%%%%%%%%%%%%%%%%%%%%%%%%%%%
%%%             Frontmatter                      %%%
%%%%%%%%%%%%%%%%%%%%%%%%%%%%%%%%%%%%%%%%%%%%%%%%%%%%
%\includeonly{chapters/4-Mala}

\begin{document}
\maketitle
\frontmatter
% %% uncomment if you have preface and/or acknowledgements

\currentpdfbookmark{Contents}{name} % adds a PDF bookmark
{\sloppy\tableofcontents}

%% sections without authors
% \include{chapters/preface}
% \addchap{\lsAcknowledgementTitle} 


\mainmatter

%%%%%%%%%%%%%%%%%%%%%%%%%%%%%%%%%%%%%%%%%%%%%%%%%%%%
%%%             Chapters                         %%%
%%%%%%%%%%%%%%%%%%%%%%%%%%%%%%%%%%%%%%%%%%%%%%%%%%%%
%
\renewcommand{\lsChapterFooterSize}{\scriptsize}
\includepaper{chapters/0-Intro}
\renewcommand{\lsChapterFooterSize}{\footnotesize}
\includepaper{chapters/1-Gast}
\renewcommand{\lsChapterFooterSize}{\scriptsize}
\includepaper{chapters/2-Saeedi}
\renewcommand{\lsChapterFooterSize}{\footnotesize}
\includepaper{chapters/3-Resi}
\includepaper{chapters/4-Mala}
\includepaper{chapters/5-Cuchet}
\includepaper{chapters/6-Giampieri}
\includepaper{chapters/7-Ray}
\includepaper{chapters/8-Kamenicka}
\renewcommand{\lsChapterFooterSize}{\scriptsize}
\includepaper{chapters/9-Kloudova}
% copy the lines above and adapt as necessary

%%%%%%%%%%%%%%%%%%%%%%%%%%%%%%%%%%%%%%%%%%%%%%%%%%%%
%%%             Backmatter                       %%%
%%%%%%%%%%%%%%%%%%%%%%%%%%%%%%%%%%%%%%%%%%%%%%%%%%%%
% There is normally no need to change the backmatter section
\sloppy 
\phantomsection
\addcontentsline{toc}{chapter}{Index} 
\addcontentsline{toc}{section}{\lsNameIndexTitle}
\ohead{\lsNameIndexTitle}
\printindex 
\cleardoublepage
  
\phantomsection 
\addcontentsline{toc}{section}{\lsLanguageIndexTitle}
\ohead{\lsLanguageIndexTitle} 
\printindex[lan] 
\cleardoublepage
  
\phantomsection 
\addcontentsline{toc}{section}{\lsSubjectIndexTitle}
\ohead{\lsSubjectIndexTitle} 
\printindex[sbj]
\ohead{} 

\end{document}
